\section{Variables, constants, and parameters}

Mathematical notation uses symbols. These symbols may play different roles. Some symbols represent variables, some represent constants, and some represent parameters.

\begin{center}
\begin{tabular}{@{}lll@{}}
\toprule
Role & Meaning & Example \\
\midrule
Variable & allowed to change inside one problem & \(x\) in \(f(x)=x^2+1\) \\
Constant & fixed value & \(3\) and \(2\) in \(f(x)=3x+2\) \\
Parameter & fixed in one case, variable between cases & \(a,b\) in \(f(x)=ax+b\) \\
\bottomrule
\end{tabular}
\end{center}

A variable is a quantity that is allowed to change. For example, in
\[
f(x)=x^2+1,
\]
the symbol \(x\) is the variable. Different values of \(x\) give different values of \(f(x)\).

A parameter is especially important because it describes a family of possibilities. The formula
\[
f(x)=ax+b
\]
describes a whole family of linear functions. For one particular function, \(a\) and \(b\) are fixed. If we choose different values of \(a\) and \(b\), we get a different function.

\begin{examplebox}
If \(a=2\) and \(b=5\), then
\[
f(x)=2x+5.
\]
If \(a=-1\) and \(b=0\), then
\[
f(x)=-x.
\]
The formula is the same, but the parameters select different members of the family.
\end{examplebox}

When reading a formula, ask which symbols change, which symbols remain fixed, and which symbols select the case being studied.
